\documentclass[userip]{lecture-digital}

\usepackage[cpalette]{sopra-listings}
\usepackage{newfloat,lecture-shortcuts}
\addto\captionsngerman{\renewcommand{\equationautorefname}{Formel}}
\addto\captionsenglish{\renewcommand{\equationautorefname}{Equation}}

\DeclareFloatingEnvironment[% 
    fileext=loc,listname={Alle Listings},name=Lst,%
    placement=tbh,within=chapter,%
]{code}
%
\usepackage{blindtext}

\titleimage{\includegraphics[width=0.6\linewidth]{exampletitleimage}}

\title{Die Kunst der alten Schule}
\subtitle{Klaus Rüdiger}
\addAuthor{Florian Sihler (florian.sihler@uni-ulm.de)}
\brief{\nohyper Eine Mitschrift von: \typesetAuthor\\{\normalsize\timestamp~$\circ$~\mail{florian.sihler@uni-ulm.de}}}
\emblem[university ulm]{Praktische Informatik}{\faCode}

\begin{document}

\frontmatter
\maketitle

\RipPage{}%[\sffamily Dokumentvariante: \LayoutName\\]

\tableofcontents
\mainmatter

\chapter{Superlangerkapitelname der gewiss nicht passen wird, denn das ist der Plan!}
\csummary{Wir beschäftigen uns mit Käse, insbesondere mit gutem, gereiften Käse! 
          Das finde ich so unglaublich subber-glasse, das ist eine zweite Zeile darüber wert, was wir in diesem Kapitel behandeln werden!
          Es wird bestimmt auch Fondue erwähnt.}

\section{Mist, er hat doch draufgepasst, das ist jetzt aber doof auch, da muss ich ihn noch länger machen}

\Blindtext[1]
\begin{figure}
    \centering\begin{tikzpicture}
        \begin{scope}
            \draw[clip] (0,0) circle [radius=1cm];
            \fill[tikzA] (1,0) circle [radius=1cm];
        \end{scope}
        \draw (1,0) circle [radius=1cm];
    \end{tikzpicture}
    \caption{Technisch superduper Bildbeschreibung die bestimmt jeder ganz ausführlich lesen wird! Die Markierung ist in \tikzAname!}
    \label{fig:example-intersect}
\end{figure}
\Blindtext[2]
Wir beziehen uns auf \autoref{fig:example-intersect}. Betrachten wir aber dazu diese tolle Liste die sich aus den Punkte \begin{inlist}
    \item Der erste Punkt 
    \item Der zweite Punkt
    \item Der Untergrund
\end{inlist}
\Blindtext[2]
\begin{table}
    \centering%
    \captionabove{A subbr Tabelle sag isch dir\ldots}
    \label{tbl:subbr-table}
    \begin{tabular}{llr}
        \toprule
            Isch & bin & Häddaaar!!! \\
        \midrule
            Öhm & isch & mein \\
            isch & bin & häggar!\\
            odder & so & glaub \\
            i \\
        \bottomrule
    \end{tabular}
\end{table}
Wir mögen den \autoref{tbl:subbr-table} fei foll ge!
\Blindtext[1]
\begin{figure*}
    \centering\begin{tikzpicture}
        \node[scale=6] {Isch bin a Bild!};
    \end{tikzpicture}
    \caption[Des isch a breites Bild!]{Des isch a Bild!}
    \label{fig:des-isch-a-bild}
\end{figure*}
Mir betrachtet jetzt \autoref{fig:des-isch-a-bild}.
\Blindtext[1]
\begin{wide}
    \begin{align}
        \ceil{\frac{2}{8}} \cdot \floor{\frac{1}{9}}& = \limn\abs{\frac{1}{n}} \label{eq:prepie}\\
        \intertext{Superduper Formulae in \autoref{eq:prepie}}
        12 \pi &= \text{Superduper}\pi
    \end{align}
\end{wide}
\Blindtext[8]
\section{Alabambi}
Ich brauche eine wirklich schöne Fußnote bitte\footnote{Kurze Fußnote.}.
% i won't use the listings features as they are much harder to control (e.g. centering)
\begin{code*}
\centering\begin{minipage}{0.8\linewidth}
\begin{java}
public static void main(String[] args) {
    System.out.println("Hallo Welt!" + 42);
}
\end{java}
\end{minipage}
\caption{\say{Hallo Welt}-Beispiel}
\label{lst:hello-world}
\end{code*}
Wir referenzieren \autoref{lst:hello-world}.
\Blindtext[6]
Hier wäre auch noch eine besonders lange Fußnote\footnote{\blindtext[1]}.

\Blinddocument
\foreach\i in {1,...,5}{
    \Blinddocument
}

\end{document}